\begin{table}[p]
\centering
\caption{Baseline Balance: Student and Class Characteristics}
\label{tab:baseline}
\begin{threeparttable}
\begin{tabular}{lrrrrr}
\toprule
 & Control Mean & Treat Mean & Std Diff & p-value & N \\
\midrule
Boy (indicator) & -- & -- & -- & -- & -- \\
Baseline EF index & -- & -- & -- & -- & -- \\
Parental education (years) & -- & -- & -- & -- & -- \\
Class size & -- & -- & -- & -- & -- \\
% Add additional baseline covariates as needed
\bottomrule
\end{tabular}
\begin{tablenotes}[flushleft]
\footnotesize
\item Notes: Placeholder table for baseline balance at randomization. Replace “--” with estimates from pre-treatment data. Report p-values from regressions of each baseline covariate on treatment with stratification fixed effects; cluster standard errors at the class level. Report N at the student level. Include counts of schools/classes (clusters) and strata sizes in the text or an accompanying panel. Apply standardized differences for comparability across measures.
\end{tablenotes}
\end{threeparttable}
\end{table}
